% Source: https://artofproblemsolving.com/community/c3h1418760_construction_geometry
% Stable identity: aops-topic-1418760-q01
\textbf{Analysis.} Extend $AC$ beyond $C$ to the point $D$ with $CD=d$. Then $AD=AC+CD=b+(c-b)=c=AB$, so $A$ is equidistant from $B$ and $D$: it lies on the perpendicular bisector of $BD$.

\textbf{Construction.}
\begin{enumerate}
\item Draw a right angle with vertex $C$ and mark $B$ on one arm with $CB=a$.
\item On the extension of the other arm beyond $C$ mark $D$ with $CD=d$.
\item Draw the perpendicular bisector of $BD$ (red); it meets the other arm at $A$.
\end{enumerate}

\begin{asy}
usepackage("amsmath");
size(5.5cm);
defaultpen(fontsize(10pt));
path ra(pair P,pair u,pair v,real s=0.2){return P+s*unit(u)--P+s*unit(u)+s*unit(v)--P+s*unit(v);}
pair ft(pair P,pair U,pair V){return U+dot(P-U,V-U)/dot(V-U,V-U)*(V-U);}
pair C=(0,0),B=(4,0),A=(0,3),D=(0,-2),M=(B+D)/2;
draw(A--B--C--cycle,linewidth(0.9)); draw(ra(C,B-C,A-C,0.25));
draw(C--D,blue+dashed); draw(B--D,blue+dashed);
draw(M+1.6*(1,-2)/sqrt(5)--A+0.6*(-1,2)/sqrt(5),red);
dot(A);dot(B);dot(C);dot(D);dot(M);
label("$A$",A,NW); label("$B$",B,E); label("$C$",C,W); label("$D$",D,SW); label("$M$",M,SE);
label("$a$",(B+C)/2,N); label("$c-b$",(C+D)/2,W,blue); label("$c$",(A+B)/2,NE);
label("\mbox{}",(min(currentpicture,user=true)+max(currentpicture,user=true))/2,black); // keeps the text after the figure black when LaTeX typesets the labels inline
\end{asy}

\textbf{Proof.} $A$ is on the perpendicular bisector, so $AB=AD=AC+d$, i.e.\ $c-b=d$; and $\angle C=90^\circ$, $BC=a$ by construction.

\textbf{When it works.} Since $c^2-b^2=a^2$, we have $c+b=\dfrac{a^2}{d}$, so $b=\dfrac{1}{2}\left(\dfrac{a^2}{d}-d\right)$, which is positive exactly when $d<a$. So a solution exists (and is unique up to congruence) if and only if $0<d<a$; this is also when the perpendicular bisector meets the arm beyond $C$. (The figure shows $a=4$, $d=2$, giving $b=3$, $c=5$.)
